%% ==========================================================================
%%  SECTION 3 --- A TAXONOMY OF PRIVACY-PRESERVING RECOMMENDATION
%%
%%  THE SPINE OF THE SURVEY.  Absorbs the superseded sections/threat_models.tex.
%%
%%  Job: give the reader axes along which to place any system in the field,
%%  then place every surveyed system on them.  This is what makes the document
%%  a survey rather than an annotated bibliography, and it is where the
%%  full-field coverage lives -- federated, DP and TEE approaches are
%%  CLASSIFIED here even though only the cryptographic line is treated at
%%  mechanism level later.
%%
%%  EVIDENCE: placement of a Tier C paper in Table~\ref{tab:taxonomy} is
%%  permitted; prose claims about Tier C papers are not.
%% ==========================================================================
\section{A Taxonomy of Approaches}
\label{sec:taxonomy}

\todo{Open by stating the problem this section solves: papers in this field are
      routinely compared as though they were alternatives when they in fact
      protect different stages under different trust assumptions with
      incomparable guarantees. The four axes below are what make the
      comparison honest.}

\subsection{Axis 1: the trust model}
\label{sec:taxonomy:trust}

\todo{Who must be trusted not to collude?
      - SINGLE SERVER (no non-collusion assumption; needs HE or single-server
        PIR --- computationally expensive)
      - TWO OR MORE NON-COLLUDING SERVERS (cheap information-theoretic or
        FSS-based protocols, at the price of an assumption about the world)
      - USER-SIDE / FEDERATED (no server holds raw data, but intermediate
        updates leak)
      - TRUSTED HARDWARE (assumes the enclave; side channels out of scope here)
      State plainly that these are not orderable by strength --- they are
      different assumptions, not more or less of one.}

\subsection{Axis 2: the privacy notion}
\label{sec:taxonomy:notion}

\todo{What does ``private'' formally mean in each case?
      - CRYPTOGRAPHIC (simulation-based; the real/ideal paradigm --- cite
        lindell2017, evans2018). Hides everything but the agreed output.
      - DIFFERENTIAL PRIVACY (cite dwork2014; mcsherry2009 applied it to
        recommendation). Bounds an individual's influence on the output;
        says nothing about who sees the computation.
      - ANONYMITY / MIXING (cite ishai2006 for the formal connection between
        anonymity and cryptographic strength).
      These compose but do not substitute. A system can be cryptographically
      private and still leak through its output --- which is exactly
      calandrino2011's attack, and the point section 9 returns to.}

\subsection{Axis 3: which pipeline stage is protected}
\label{sec:taxonomy:stage}

\todo{Map onto Figure~\ref{fig:pipeline}: collection / training / serving /
      delivery. THIS IS THE AXIS THE FIELD IS LEAST EXPLICIT ABOUT, and the
      one that makes the survey's contribution. Most work protects training;
      considerably less protects delivery; very little protects both.}

\subsection{Axis 4: the adversary}
\label{sec:taxonomy:adversary}

\todo{Semi-honest versus malicious; static versus adaptive corruption;
      the collusion threshold. Note that almost all deployed-scale systems in
      this survey are semi-honest, and that this is a substantive limitation
      rather than a technicality.}

\begin{table}[t]
\centering
\todo{THE MASTER CLASSIFICATION TABLE --- the single most valuable object in
      this survey. Rows: every surveyed system. Columns: trust model
      (axis 1) | privacy notion (axis 2) | pipeline stage(s) covered
      (axis 3) | adversary (axis 4) | evaluated scale.
      Tier C papers may appear here, since placement follows from title and
      abstract. Sort by pipeline stage so the sparsity of the delivery column
      is visible at a glance --- that sparsity is a finding, and it should be
      readable from the table without prose.
      Use booktabs; this will need \small or a sideways float.}
\caption{Classification of surveyed systems along the four axes of
         Section~\ref{sec:taxonomy}.}
\label{tab:taxonomy}
\end{table}

\subsection{Reading the table}
\label{sec:taxonomy:reading}

\todo{Two or three paragraphs drawing out what the table shows --- and ONLY
      what it shows. Expected observations, to be confirmed once the table is
      populated rather than assumed now:
      (i) the training column is dense and the delivery column sparse;
      (ii) the cheapest protocols consistently buy their speed with a
           non-collusion assumption;
      (iii) systems at the largest evaluated scale are semi-honest.
      \verify{Do not write these until the table is filled --- they are
      predictions, and the table may not support them.}}
